\section{Phase 2: Implementation and Latency Evaluation}\label{sec:impl}

\subsection{Implementation}
Each component of Fig.~\ref{fig:arch} runs in its own container, defined in one
Compose file: Redis, the server and a client container from which we run the
command-line client, the load generator and the tests. The server and client
images are built from one Dockerfile based on \texttt{python:3.12-slim}; the
\texttt{python:3.9-slim-buster} image from the assignment no longer builds because
Debian Buster is end-of-life. The corpus consists of eight books from Project
Gutenberg (\textbf{[N]} words). The server splits every text into lowercase words
once at startup; a cache miss counts the keyword in that word list. Every request
uses a new connection, because the load balancer of Phase~3 can only choose a
server when a connection is opened.

\subsection{Experiments}
Latency is measured at the client, from opening the connection to receiving the
reply. The load generator sends a fixed list of 5{,}000 random requests (random
text, random keyword from the 500 most common words of a book without the 50 most
common ones, fixed seed) at a fixed rate. Each request is sent from its own thread
at its scheduled time, whether or not earlier replies have arrived, so slow replies
do not lower the request rate. We used the rates \textbf{[R1--R5]}\,req/s and ran
each rate three times. Redis is emptied before every run, so every run starts with
an empty cache and sees the same sequence of hits and misses.

\begin{table}[t]
\caption{Phase 2 experiment setup}
\label{tab:setup}
\centering
\footnotesize
\begin{tabular}{@{}p{0.25\columnwidth}p{0.69\columnwidth}@{}}
\hline
Host & \textbf{[CPU, cores, RAM, OS]}; Docker \textbf{[version]} \\
Software & Python 3.12, RPyC 6.0.2, redis-py 8.1.0, Redis 7.4 \\
Corpus & 8 Project Gutenberg books, \textbf{[N]} words \\
Requests & 5{,}000 per run, random text and keyword, fixed seed \\
Rates & \textbf{[R1, R2, R3, R4, R5]} requests/s \\
Runs & 3 per rate, each with an empty cache \\
Metrics & mean and 99th-percentile latency \\
\hline
\end{tabular}
\end{table}

\subsection{Results}
Fig.~\ref{fig:shot-p2} shows that the server answers requests and uses the cache.
Figs.~\ref{fig:mean} and~\ref{fig:p99} show the mean and the 99th-percentile
latency per rate (average of three runs); Table~\ref{tab:latency} lists the values.
\textbf{[All requests were answered without errors.]}
% TODO after the runs: replace the bracketed values and keep only what your figures show.
At the lowest rate the mean latency is \textbf{[x]}\,ms. Most of it is a fixed
cost per request: opening a TCP connection and setting up the RPyC connection,
which a cache hit does not remove; a miss also counts the word in the text, which
takes longest for the largest book. Up to \textbf{[R]}\,req/s the mean stays almost
flat, while the 99th percentile grows from \textbf{[a]} to \textbf{[b]}\,ms. The
tail grows first because requests that arrive close together have to wait for
each other: the server handles them in threads, but Python runs only one thread at
a time, so the work is done one request after another. Mixing fast cache hits with
slow misses makes these waiting times uneven, so the tail grows well before the
server is fully busy~\cite{dean2013tail}. At \textbf{[R5]}\,req/s the 99th
percentile reaches \textbf{[c]}\,ms: one server does not scale, which motivates
replicating it behind a load balancer (Phase~3).

\begin{figure}[t]
\centering
\includegraphics[width=\columnwidth]{figures/screenshot_phase2}
\caption{The client asks for the same keyword twice (first a cache miss, then a cache
hit) and for the hot keywords; below, the matching lines of the server log.}
\label{fig:shot-p2}
\end{figure}

\begin{figure}[t]
\centering
\includegraphics[width=\columnwidth]{figures/phase2_mean}
\caption{Average latency per request rate, one server (average of 3 runs).}
\label{fig:mean}
\end{figure}

\begin{figure}[t]
\centering
\includegraphics[width=\columnwidth]{figures/phase2_p99}
\caption{99th-percentile (tail) latency per request rate, one server (average of 3 runs).}
\label{fig:p99}
\end{figure}

\begin{table}[t]
\caption{Latency per request rate, one server (ms, average of 3 runs)}
\label{tab:latency}
\centering
\footnotesize
\begin{tabular}{rrr}
\hline
Rate (req/s) & Mean & p99 \\
\hline
\textbf{[R1]} & & \\
\textbf{[R2]} & & \\
\textbf{[R3]} & & \\
\textbf{[R4]} & & \\
\textbf{[R5]} & & \\
\hline
\end{tabular}
\end{table}
